% =============================================================================
% LIMITATIONS - CITA Paper in ALKALI Format
% =============================================================================

\vspace{-2mm}
\section{Limitations}
\label{sec:limitations}

We acknowledge the current limitations of CITA and outline directions for future work.

\vspace{-2mm}
\subsection{Experimental Scope}
\label{sec:experimental_scope}

\begin{table}[H]
\centering
\small
\begin{tabularx}{\columnwidth}{@{}lX@{}}
\toprule
\textbf{Limitation} & \textbf{Impact \& Mitigation Path} \\
\midrule
Single model (Llama-3.1-8B) & Validate on GPT, Mistral, Gemma; scale to 70B \\
English-only evaluation & Extend to multilingual instruction transfer \\
10 instruction types & Add domain-specific types (medical, legal) \\
Single dataset (PKU-SafeRLHF) & Test on diverse preference datasets \\
\bottomrule
\end{tabularx}
\caption{Current experimental limitations and mitigation paths.}
\label{tab:limitations}
\vspace{-4mm}
\end{table}

\paragraph{Model Generalization.} We evaluate only Llama-3.1-8B~\cite{dubey2024llama}. Future work should validate CITA on diverse architectures~\cite{jiang2023mistral,zhang2022opt,openai2023gpt4} and scales (7B--70B) to establish generality across model families.

\paragraph{Instruction Coverage.} The 10 instruction types in ISD cover common alignment dimensions but may not capture all real-world policy requirements. Domain-specific instructions (medical, legal, educational) require dedicated evaluation.

\paragraph{Multilingual Alignment.} All experiments are English-only. Instruction-conditioned alignment in multilingual settings raises questions about cross-lingual instruction transfer and cultural alignment differences.

\vspace{-2mm}
\subsection{Deployment Considerations}
\label{sec:deployment}

\paragraph{Inference Overhead.} CITA requires alignment instructions in every prompt, increasing context length by 10--40 tokens. While minimal for modern LLMs, this may matter for latency-critical applications or extremely long contexts.

\paragraph{Instruction Specification.} Deployers must carefully craft alignment instructions for each use case. Poorly specified instructions may lead to unexpected behaviors---instruction engineering becomes a new deployment concern.

\vspace{-2mm}
\subsection{Safety \& Ethical Considerations}
\label{sec:safety_ethics}

\paragraph{Dual-Use Risk.} Instruction-conditioned alignment could be misused to make models \textit{less} safe via adversarial instructions~\cite{zou2023universal,carlini2023aligned,pace2024west}. Mitigations include:
\begin{itemize}[leftmargin=*,itemsep=1pt]
    \item Hierarchical instruction handling (system $>$ user)
    \item Instruction validation/filtering at deployment
    \item Constitutional constraints~\cite{bai2022constitutional} as non-negotiable baselines
\end{itemize}

\paragraph{Alignment Washing.} Organizations might claim ``aligned'' behavior while using permissive instructions. Transparency about instruction policies and third-party auditing are essential safeguards.

\paragraph{Control Boundaries.} CITA enables user-influenced alignment, raising questions about who controls behavioral policy. We advocate for tiered systems: users adjust within deployer-set bounds, which operate within regulatory constraints.

\vspace{-2mm}
\subsection{Future Work}
\label{sec:future_work}

\begin{enumerate}[leftmargin=*,itemsep=2pt]
    \item \textbf{Hierarchical Instructions}: Multi-level handling where system policies override user preferences

    \item \textbf{Instruction Compositionality}: Combining multiple instructions (``be concise AND professional'') meaningfully

    \item \textbf{Instruction Robustness}: Resistance to instruction injection attacks

    \item \textbf{Theoretical Analysis}: Convergence guarantees and sample complexity for instruction-conditioned learning

    \item \textbf{Multi-Modal CITA}: Extension to vision-language models with cross-modal instructions
\end{enumerate}
